\documentclass[11pt,a4paper]{article}
\usepackage[utf8]{inputenc}
\usepackage[T1]{fontenc}
\usepackage{amsmath,amssymb,amsthm}
\usepackage{geometry}
\geometry{margin=2.5cm}
\theoremstyle{plain}
\newtheorem{theorem}{Theorem}[section]
\newtheorem{proposition}[theorem]{Proposition}
\newtheorem{lemma}[theorem]{Lemma}
\theoremstyle{definition}
\newtheorem{definition}[theorem]{Definition}
\title{Soluciones Tests 36-40: Cálculo}
\author{Mathematical AI Agent}
\date{\today}
\begin{document}
\maketitle

\section{Problemas}

\subsection{Teorema de Rolle}

\begin{theorem}[Rolle]
Sea $f: [a,b] \to \mathbb{R}$ una función que verifica:
\begin{enumerate}
    \item $f$ es continua en $[a,b]$,
    \item $f$ es derivable en $(a,b)$,
    \item $f(a) = f(b)$.
\end{enumerate}
Entonces existe al menos un punto $c \in (a,b)$ tal que $f'(c) = 0$.
\end{theorem}

\begin{proof}
Sea $f: [a,b] \to \mathbb{R}$ que satisface las hipótesis del teorema.

Por el \textbf{teorema de Weierstrass}, dado que $f$ es continua en el intervalo cerrado y acotado $[a,b]$, existen puntos $m, M \in [a,b]$ tales que:
\[
f(m) = \min_{x \in [a,b]} f(x), \quad f(M) = \max_{x \in [a,b]} f(x).
\]

\textbf{Caso 1:} Si $f(m) = f(M)$, entonces $f$ es constante en $[a,b]$, pues para todo $x \in [a,b]$:
\[
f(m) \leq f(x) \leq f(M) = f(m).
\]
En este caso, $f'(x) = 0$ para todo $x \in (a,b)$, y cualquier punto $c \in (a,b)$ cumple la conclusión.

\textbf{Caso 2:} Si $f(m) \neq f(M)$, como $f(a) = f(b)$, al menos uno de los extremos $m$ o $M$ no puede ser igual a $a$ ni $b$. Supongamos, sin pérdida de generalidad, que $c \in (a,b)$ es un punto donde se alcanza el máximo o mínimo (es decir, $f(c) = f(m)$ o $f(c) = f(M)$).

Por el \textbf{teorema de Fermat}, si $f$ alcanza un extremo local en un punto interior $c \in (a,b)$ donde $f$ es derivable, entonces:
\[
f'(c) = 0.
\]

Alternativamente, podemos demostrarlo directamente usando el \textbf{criterio de crecimiento}: si $f'(c) > 0$ en algún punto, existiría una vecindad donde $f$ es creciente, contradiciendo que $c$ sea un máximo local. Análogamente, si $f'(c) < 0$, existiría una vecindad donde $f$ es decreciente, contradiciendo que $c$ sea un mínimo local.

Por lo tanto, existe $c \in (a,b)$ tal que $f'(c) = 0$.
\end{proof}

\subsection{Teorema del Valor Medio}

\begin{theorem}[Del Valor Medio]
Sea $f: [a,b] \to \mathbb{R}$ una función que verifica:
\begin{enumerate}
    \item $f$ es continua en $[a,b]$,
    \item $f$ es derivable en $(a,b)$.
\end{enumerate}
Entonces existe al menos un punto $c \in (a,b)$ tal que:
\[
f'(c) = \frac{f(b) - f(a)}{b - a}.
\]
\end{theorem}

\begin{proof}
Definimos la función auxiliar $g: [a,b] \to \mathbb{R}$ por:
\[
g(x) = f(x) - f(a) - \frac{f(b) - f(a)}{b - a}(x - a).
\]

Verificamos que $g$ satisface las hipótesis del \textbf{teorema de Rolle}:

\begin{itemize}
    \item $g$ es continua en $[a,b]$: La función $f$ es continua por hipótesis, y las funciones $x \mapsto f(a)$ y $x \mapsto \frac{f(b)-f(a)}{b-a}(x-a)$ son continuas (son funciones lineales o constantes).
    
    \item $g$ es derivable en $(a,b)$: Análogamente, $f$ es derivable por hipótesis, y las demás funciones que componen $g$ son derivables.
    
    \item Verificación de las condiciones en los extremos:
    \[
    g(a) = f(a) - f(a) - \frac{f(b) - f(a)}{b - a}(a - a) = 0,
    \]
    \[
    g(b) = f(b) - f(a) - \frac{f(b) - f(a)}{b - a}(b - a) = f(b) - f(a) - (f(b) - f(a)) = 0.
    \]
\end{itemize}

Por el \textbf{teorema de Rolle}, existe $c \in (a,b)$ tal que $g'(c) = 0$.

Calculamos $g'(x)$:
\[
g'(x) = f'(x) - \frac{f(b) - f(a)}{b - a}.
\]

Evaluando en $c$:
\[
g'(c) = f'(c) - \frac{f(b) - f(a)}{b - a} = 0.
\]

Por lo tanto:
\[
f'(c) = \frac{f(b) - f(a)}{b - a}.
\]
\end{proof}

\subsection{Derivada Nula Implica Constante}

\begin{theorem}
Sea $f: I \to \mathbb{R}$ una función derivable en un intervalo abierto $I$, y sea $f'(x) = 0$ para todo $x \in I$. Entonces $f$ es constante en $I$.
\end{theorem}

\begin{proof}
Sea $a, b \in I$ con $a < b$. El intervalo $[a,b] \subseteq I$ (pues $I$ es abierto, podemos considerar un intervalo más pequeño si es necesario, pero la propiedad se extiende por transitividad).

Aplicamos el \textbf{teorema del valor medio} a $f$ en $[a,b]$:
\[
f(b) - f(a) = f'(c)(b - a)
\]
para algún $c \in (a,b)$.

Por hipótesis, $f'(c) = 0$, por lo tanto:
\[
f(b) - f(a) = 0 \cdot (b - a) = 0,
\]
lo que implica $f(b) = f(a)$.

Dado que $a$ y $b$ son arbitrarios en $I$, concluimos que $f(x) = f(a)$ para todo $x \in I$, es decir, $f$ es constante en $I$.
\end{proof}

\subsection{Derivada Positiva Implica Creciente}

\begin{theorem}
Sea $f: I \to \mathbb{R}$ una función derivable en un intervalo abierto $I$, y sea $f'(x) > 0$ para todo $x \in I$. Entonces $f$ es estrictamente creciente en $I$.
\end{theorem}

\begin{proof}
Sea $a, b \in I$ con $a < b$.

Aplicamos el \textbf{teorema del valor medio} a $f$ en $[a,b]$:
\[
f(b) - f(a) = f'(c)(b - a)
\]
para algún $c \in (a,b)$.

Por hipótesis, $f'(c) > 0$, y como $a < b$ tenemos $b - a > 0$. Por lo tanto:
\[
f(b) - f(a) = f'(c)(b - a) > 0,
\]
lo que implica $f(b) > f(a)$.

Dado que $a$ y $b$ son arbitrarios con $a < b$, concluimos que $f$ es estrictamente creciente en $I$.
\end{proof}

\subsection{Sucesiones Monótonas y Acotadas}

\begin{theorem}
Toda sucesión monótona y acotada en $\mathbb{R}$ converge.
\end{theorem}

\begin{proof}
Sea $(x_n)_{n \in \mathbb{N}}$ una sucesión monótona y acotada en $\mathbb{R}$.

\textbf{Caso 1:} Supongamos que $(x_n)$ es creciente (es decir, $x_n \leq x_{n+1}$ para todo $n$) y acotada superiormente.

Por el \textbf{axioma de completitud} de $\mathbb{R}$ (o \textbf{propiedad del supremo}), el conjunto $A = \{x_n : n \in \mathbb{N}\}$ tiene supremo. Definimos:
\[
L = \sup_{n \in \mathbb{N}} x_n.
\]

Probaremos que $\lim_{n \to \infty} x_n = L$. Sea $\varepsilon > 0$ arbitrario.

Como $L$ es el supremo de $A$, $L - \varepsilon$ no es cota superior de $A$. Por lo tanto, existe un índice $N \in \mathbb{N}$ tal que:
\[
x_N > L - \varepsilon.
\]

Dado que $(x_n)$ es creciente, para todo $n \geq N$:
\[
x_n \geq x_N > L - \varepsilon.
\]

Además, como $L$ es cota superior:
\[
x_n \leq L < L + \varepsilon.
\]

Por lo tanto, para todo $n \geq N$:
\[
|x_n - L| = L - x_n < \varepsilon,
\]
ya que $x_n \leq L$. Esto demuestra que $\lim_{n \to \infty} x_n = L$.

\textbf{Caso 2:} Si $(x_n)$ es decreciente y acotada inferiormente, definimos $L = \inf_{n \in \mathbb{N}} x_n$ y la demostración es análoga.

Por lo tanto, toda sucesión monótona y acotada en $\mathbb{R}$ converge.
\end{proof}

\end{document}